\section*{Bab \ref{ch:g4:mult} --- Mengalikan Bilangan yang Lebih Besar}

\begin{solution}{exo:g4:mult:1}
$470$; \quad $2\,300$; \quad $720$; \quad $4\,500$; \quad
$2\,100$.
\end{solution}

\begin{solution}{exo:g4:mult:2}
$34 \times 21 = 34 \times 20 + 34 = 680 + 34 = 714$; cara bersusun
memberi baris yang sama: $34$ dan $680$.
\end{solution}

\begin{solution}{exo:g4:mult:3}
$63 \times 24 = 1\,512$; \quad
$85 \times 47 = 3\,995$; \quad
$126 \times 53 = 6\,678$.
\end{solution}

\begin{solution}{exo:g4:mult:4}
$208 \times 34$: taksiran $200 \times 34 = 6\,800$; tepatnya
$7\,072$.

$517 \times 62$: taksiran $500 \times 60 = 30\,000$; tepatnya
$32\,054$.
\end{solution}

\begin{solution}{exo:g4:mult:5}
$45 \times 11 = 450 + 45 = 495$.

$32 \times 99 = 3\,200 - 32 = 3\,168$.

$24 \times 45 = 12 \times 90 = 1\,080$ (atau sekali lagi:
$6 \times 180 = 1\,080$).
\end{solution}

\begin{solution}{exo:g4:mult:6}
$26 \times 18 = 468$ kursi.
\end{solution}

\begin{solution}{exo:g4:mult:7}
Taksiran: $30 \times 50 = 1\,500$. Tepatnya:
$32 \times 48 = 1\,536$ peti.
\end{solution}

\begin{solution}{exo:g4:mult:8}
Ia lupa bahwa angka $2$ pada $23$ menghitung \emph{puluhan}: baris
kedua seharusnya $57 \times 20 = 1\,140$, bukan $114$. \omterm{def:g1:addition:def}{Jumlah} yang
benar: $171 + 1\,140 = 1\,311$.
\end{solution}

\begin{solution}{exo:g4:mult:9}
Bunga tulip: $14 \times 35 = 490$. Bunga bakung: $12 \times 28 = 336$.
Seluruhnya: $490 + 336 = 826$ bunga.
\end{solution}

\begin{solution}{exo:g4:mult:10}
Taksiran: $70 \times 30 = 2\,100$ kg. Tepatnya:
$67 \times 30 = 2\,010$ kg.
\end{solution}

\begin{solution}{exo:g4:mult:11}
Persegi panjang $15 \times 12$ yang tersusun dari persegi satuan
dapat dipotong oleh sebuah garis tegak menjadi bagian
$15 \times 10$ dan bagian $15 \times 2$: membilang tiap bagian
memberi $150 + 30 = 180$. Potongan yang lain berhasil sama baiknya,
misalnya $15 = 10 + 5$: $10 \times 12 + 5 \times 12 = 120 + 60
= 180$. Cara memotong mana pun tidak mengubah banyaknya persegi
seluruhnya.

\end{solution}
